% Generated draft: requires mathematical and source review.
\documentclass[11pt]{article}
\usepackage[margin=25mm]{geometry}
\usepackage{fontspec}
\setmainfont{Latin Modern Roman}
\usepackage{amsmath,amssymb,mathtools,unicode-math}
\setmathfont{Latin Modern Math}
\usepackage{enumitem,needspace,xurl,hyperref}
\setlength{\parindent}{0pt}
\setlength{\parskip}{5pt}
\setlength{\emergencystretch}{3em}
\newcommand{\N}{\mathbb{N}}
\newcommand{\Z}{\mathbb{Z}}
\newcommand{\Q}{\mathbb{Q}}
\newcommand{\R}{\mathbb{R}}
\newcommand{\C}{\mathbb{C}}
\newcommand{\F}{\mathbb{F}}
\newcommand{\A}{\mathbb{A}}
\newcommand{\PP}{\mathbb{P}}
\newcommand{\E}{\mathbb{E}}
\newcommand{\eps}{\varepsilon}
\newcommand{\dd}{\,\mathrm{d}}
\begin{document}
\section*{UnsolvedMath ID: 30006223}
\textbf{Source problem code:} catalogue entry
\subsection{Definitions and mathematical statement}
Write
\[
D^3=\{x\in\R^3:\|x\|\leq 1\},\qquad S^1=\{x\in\R^2:\|x\|=1\}.
\]
Thus \(D^3\times S^1\) is a compact four-dimensional manifold with boundary \(S^2\times S^1\). The notation \(\deg_{\mathrm{exp}}(X)\) denotes the expander degree referred to in the supplied record and its Bamler--Chen reference. The record does not give the definition of this invariant, its normalization, its admissible class of spaces, or any auxiliary data required to define it.

In the usual smooth manifold convention, the boundary connected sum \(X_1\mathbin{\#_\partial}X_2\) of compact smooth \(n\)-manifolds with nonempty boundary is formed by attaching a one-handle \(D^1\times D^{n-1}\) to their disjoint union, with one attaching disk in each boundary, and smoothing corners. This explains the notation only: the record does not specify whether this category, additional structures, or restrictions on the attaching data are intended.

\textbf{Requested relation and value.} For every pair \(X_1,X_2\) satisfying the hypotheses of the intended expander-degree theory, and for which the boundary connected sum is admissible, determine a general relation among
\[
\deg_{\mathrm{exp}}(X_1),\qquad
\deg_{\mathrm{exp}}(X_2),\qquad
\deg_{\mathrm{exp}}(X_1\mathbin{\#_\partial}X_2).
\]
In particular, determine
\[
\deg_{\mathrm{exp}}(D^3\times S^1).
\]
The source asks whether there is any relation; it does not propose a specific identity. The admissibility hypotheses, degree definition, and dependence on auxiliary or attaching data are unresolved in the supplied record, so this is a conditional formulation of the requested problem rather than a fully specified standalone statement.
\subsection{Short English statement}
Determine how expander degree behaves under boundary connected sum, and compute its value on \(D^3\times S^1\). The supplied record does not define expander degree or specify the hypotheses on the spaces, so these must be recovered from the cited theory before the question has a complete standalone formulation.
\subsection{Sources}
\begin{itemize}
\item Local UnsolvedMath registry, supplied record with permanent ID \(30006223\), title \emph{Expander Degree Under Boundary Connected Sums}, source problem code \texttt{catalogue entry}, in the Oberwolfach Reports Open Problems collection. The project registry paths are \texttt{data/unsolvedmath/problems.json} and the supplementary \texttt{catalog/problems.json}; no independent inspection of either file is claimed.
\item \emph{MATRIX-MFO Tandem Workshop: Nonlinear Geometric Diffusion Equations} (2025), Oberwolfach Reports, report 11/2025, pp.~485--530, DOI \url{https://doi.org/10.4171/owr/2025/11}. The supplied background describes Oberwolfach Reports as documenting mathematical workshops and their open-problem discussions. No individual proposer or precise location within the report is supplied.
\item Bamler--Chen, \emph{Expander degree theory}, \url{https://doi.org/10.1002/cpa.70024}. This is a supplied literature reference; its contents have not been independently consulted here.
\item Supplied literature assessment, dated exactly \texttt{2026-08-22}: the problem is labelled open, and the assessment states that the expander-degree paper proves the degree for \(D^4\) and several orbifolds while listing the value for \(S^1\times D^3\) and boundary-connected-sum behavior as open questions. It consequently describes the requested \(D^3\times S^1\) value as unresolved. This is supplied dated curation derived from the OWR record, not an independently verified literature assessment; its status and cited claims require independent verification.
\end{itemize}
\end{document}
